\documentclass[10pt]{article}

\usepackage[utf8]{inputenc}

\usepackage[T1]{fontenc}

\usepackage{amsmath}

\usepackage{amsfonts}

\usepackage{amssymb}

\usepackage[version=4]{mhchem}

\usepackage{stmaryrd}

\usepackage{bbold}

\usepackage{graphicx}

\usepackage[export]{adjustbox}

\graphicspath{ {./images/} }



\begin{document}

ОКА ТЕОРЕМЫ - теоремы о классич. проблемах теории функций многих комплексных переменных, впервые доказанные К. Ока в 1930-50 (см. [1]).



\begin{enumerate}

  \item О. т. о К узена проблемах: первая проблема Кузева разрешима в любой области голоморфности в $\mathbb{C}^{n} ;$ вторая проблема Кузена разрешима в любой области голоморфности $D \subset \mathbb{C}^{n}$, гомеоморфной $D_{1} \times \ldots \times D_{n}$, где все области $D_{v} \subset \mathbb{C}$, кроме, возможно, одной, односвязны.



  \item О. т. о Леви проблеме: всякая псевдовыпуклая риманова область является областью голоморфности.



\end{enumerate}



Первоначально эта теорема была доказана К. Ока для размерности $n=2$; в случае произвольной размерности она доказана К. Ока и др. математиками.



\begin{enumerate}

  \setcounter{enumi}{2}

  \item О к а - В е йл я т е о р е ма: пусть $D-$ область в $\mathbb{C}^{n}$ и компакт $K \subset D$ совпадает со своей оболочкой относительно алгебры $G(D)$ всех голоморфных в $D$ функций; тогда для любой функции $f$, голоморфной в окрестности $K$, і любого $\varepsilon>0$ найдется функция $F \in G(D)$ такая, что

\end{enumerate}



\[

\max _{\mathrm{K}}|f-F|<\varepsilon .

\]



Эта фундаментальная теорема теории голоморфных приближениї широко применяется в комплексном и функциональном анализе.



\begin{enumerate}

  \setcounter{enumi}{3}

  \item О. т. о к о ге ре н т н о с т п: пусть $6-$ пучок голоморфных функций на комплексном многообразии $X$; тогда для любого натурального числа $p$ любої локально конечно порожденныйі подпучок пучка $6^{p}=6 \times \ldots \times 6$ ( $p$ раз) является когерентным аналитическим пучко.и.

\end{enumerate}



Это одна из основных теорем т. н. теории Ока Картана, к-рая существенно используется при доказательстве Картана теорем А и В.



Jum.: [1] О К a K., Sur les fonctions analitiques plusieurs variables, Tokyo, 1961; '2] X e p м а н д е p Л., Введение в теорию функций нескольких компјексных переменных, пер. с англ., М., 1968; [3] Г а н н и н г Р., Р о с с и Х., Аналитические функции многих комилексных персменных, пер. с англ., М., 1969.



\includegraphics[max width=\textwidth, center]{2023_07_08_0577701e8d7fa9b6422fg-1}

п а ктном р а с ши и р н ии bX- конечное семейство $\left\{U_{1}, \ldots, U_{k}\right\}$ открытых в $X$ множеств такое, что множество $K \backslash U_{1} \cup \ldots \cup U_{k}$ бикомпактно и $b X=$ $=K \cup \widetilde{U}_{1} \cup \ldots \cup \tilde{U}_{k}$, где $\tilde{U}_{k}-$ наибольшее открытое в $b X$ множество, высекающее на $X$ множество $U_{i}(X$ предполагается вполне регулярным). Понятие О.п. $X$ в $b X$ совпадает с понятием близостного продолжаемого окаймления пространства близости $X$ (близость на $X$ индуцирована расширением $b X$ ), формулируемом в близостных терминах: кроме бикоміактности $K$ требуется, чтобы для любой окрестности $\widehat{\sigma}_{K}$ семейство $\left\{6_{K}, U_{1}, \ldots, U_{k}\right\}$ было равномерным покрытием пространства $X$. Просто о к а й м л е н и е м пространства $X$ наз. его окаймление в бикомпактном расширении Стоуна - Чеха. На языке окаймлений формулируется ряд теорем о размерности наростов бикомпактных расширений топологических и близостных пространств. Јum.: [1] С м и р н о в Ю. М., "Матем. сб.», 1966, т. 71, № 4,

с. $454-82$.

В. Федориук.



ОКАЙМЛЕНИЯ МЕТОД - метод решения системы линеӥных алгебраич. уравнений $A x=b$ с невырожденвой матрицей, обращения матрицы и вычисления определителя, основанный на рекурревтном переходе от решения задачи с матрицей



\[

A_{k-1}=\left\|\begin{array}{lll}

a_{1,1} & \ldots & a_{1, k-1} \\

\ldots & \ldots & \ldots \\

a_{k-1,1} & \ldots & a_{k-1, k-1}

\end{array}\right\|

\]



к решению задачи с матрицей $A_{k}$, рассматриваемой как результат окаймления $A_{k-1}$.



Вычислительная схема О. м. для обращения матриц такова. Пусть $A_{k-1}$ - невырожденная матрица. Для обращения матрицы $A_{k}$ используется представление



\[

A_{k}=\left\|\begin{array}{cc}

A_{k-1} & u_{k} \\

v_{k} & a_{k k}

\end{array}\right\|

\]



где



$u_{k}=\left(a_{1, k}, a_{2, k}, \ldots, a_{k+1, k}\right)^{T}, \quad v_{k}=\left(a_{k, 1}, a_{k, 2}, \ldots, a_{k, k-1}\right)$.



Тогда



\[

A_{\bar{k}}^{-1}=\| \begin{array}{cc}

A_{k-1}^{-1}+\frac{A_{k-1}^{-1} u_{k} v_{k} A_{k-1}^{-1}}{\alpha_{k}} & -\frac{A_{k-1}^{-1} u_{k}}{\alpha_{k}} \\

-\frac{v_{k} A_{k-1}^{-1}}{\alpha_{k}} & \frac{1}{\alpha_{k}} \\

\alpha_{k}=a_{k k}-v_{k} A_{k-1}^{-1} u_{k} .

\end{array}

\]



Последовательное обращение матриц $A_{1}, A_{2}, \ldots, A_{n}$ по этой схеме дает матрицу $A^{-1}$.



Описанная схема О. м. пригодна лишь для матриц с отличными от нуля главными минорами. В общем случае следует принять схему О. м. с выбором главного әлемснта. В этой схеме в качестве окаймляющих строки и столбца берутся те, для к-рых $\alpha_{k}=a_{k k}-v_{k} A_{k-1}^{-1} u_{k}$ будет максимальным по модулю. Тогда вычисленная матрица будет отличаться от $A^{-1}$ лишь перестановкой строк и столбцов [1].



По быстродействию О. м. не уступает самым быстродействующим из прямых методов обращения матрицы.



О. м. позволяет әффективно обращать треугольные матрицы. Если $A_{k-1}$ - правая треугольная матрица, то в (1)



\[

v_{k}=0 \text { и } A_{k}^{-1}=\| \begin{array}{cc}

A_{k-1}^{-1} & -\frac{A_{\bar{k}-1}^{-1} u_{k}}{a_{k k}} \\

0 & \frac{1}{a_{k k}}

\end{array} .

\]



Объем вычислений в этом случае уменьшается в 6 раз.



Особенно әффективен О. м. при обращении эрмитовых положительно определенных матриц. Для этих матриц не нужно применять схему выбора главного әлемента. Кроме того, они могут быть заданы лишь половиной своих элементов. Вычислительная схема в этом случае упрощается:



\[

\begin{gathered}

A_{k}^{-1}=\| \begin{array}{rr}

A_{k-1}^{-1}+\frac{P_{k} P_{k}^{*}}{\alpha_{k}} & -\frac{P_{k}}{\alpha_{k}} \\

-\frac{P_{k}^{*}}{\alpha_{k}} & \frac{1}{\alpha_{k}}

\end{array},, \\

P_{k}=A_{k-1}^{-1} u_{k}, \alpha_{k}=a_{k k}-u_{k}^{*} P_{k} .

\end{gathered}

\]





\end{document}